% 44.5 — 新 SOP 重画（2026-10-20）
% 真值 = canvas.png（1135x459，1pt=1px，y 向下）
% 左：刻度梳；中：弧形箭头 + f3；右：4x4 方格 + 折线图案
\documentclass[12pt]{standalone}
\usepackage{fontspec}
\setsansfont{TeX Gyre Heros}
\usepackage{tikz}
\usetikzlibrary{arrows.meta}
% =====================================================================
% 统一符号层（2026-10-20）
% 用法：在 fig_<号>.tex 的 preamble 里 % =====================================================================
% 统一符号层（2026-10-20）
% 用法：在 fig_<号>.tex 的 preamble 里 % =====================================================================
% 统一符号层（2026-10-20）
% 用法：在 fig_<号>.tex 的 preamble 里 \input{../_symbols/symbols.tex}
% 依赖：fontspec（已 \setsansfont{TeX Gyre Heros}）、graphicx 不需要
% 规则见同目录《符号验收清单.md》
% =====================================================================

% ---- 字体角色（全部走 sans = TeX Gyre Heros）----
\newcommand{\SIt}[1]{\textsf{\textit{#1}}}          % 斜体变量
\newcommand{\SBf}[1]{\textsf{\textbf{#1}}}          % 粗体
\newcommand{\SBfIt}[1]{\textsf{\textbf{\textit{#1}}}} % 粗斜体变量
\newcommand{\SNum}[1]{\textsf{\textup{#1}}}         % 直立数字

% ---- 上下标：字号/位置写死，避免数学样式与手调混用 ----
% 审计结论（2026-10-20）：源图下标/底字高比在 0.77–0.92em，0.72em 偏小。
% 默认 0.80em；仍需按源图逐图标定（44.4/44.5 约 0.85–0.92，39/41/46 约 0.77–0.80）。
\newcommand{\FSup}[2]{#1\raisebox{0.30em}{\fontsize{0.72em}{0.82em}\selectfont #2}}
\newcommand{\FSub}[2]{#1\raisebox{-0.30em}{\fontsize{0.80em}{0.92em}\selectfont #2}}

% ---- 分数：数字直立、方头粗杆；杆宽 = max(分子宽,分母宽) + 2pt ----
% 杆厚 4pt（配合 20–26pt 标签）；小字号图可局部调 0.17em
% 2026-10-20 修正：原 \shortstack 负 kern 不随字号缩放（23pt 下总高仅 30pt、杆贴字），
% 改为 offinterlineskip + em 定间距：上间隙 0.12em、下间隙 0.13em
% （37.2 源图 23pt 实测四标签上/下间隙均值 2.6/4.0pt；取近平衡且总高保持 42–43pt）。
\newcommand{\FFrac}[2]{%
  \begingroup
  \setbox0=\hbox{#1}\setbox2=\hbox{#2}%
  \dimen0=\wd0 \ifdim\wd2>\dimen0 \dimen0=\wd2 \fi
  \advance\dimen0 by 2pt
  \leavevmode\vbox{\offinterlineskip
    \hbox to \dimen0{\hss\box0\hss}%
    \kern0.12em
    \hbox{\rule{\dimen0}{4pt}}%
    \kern0.13em
    \hbox to \dimen0{\hss\box2\hss}%
  }%
  \endgroup}

% ---- 真减号（审计：源图减号 15–18px 宽、3–5px 厚；Heros 常规字重偏细）----
% 可选参数 = 宽度（em）：正文 0.58；上标内用 \FMinus[0.72]
\newcommand{\FMinus}[1][0.58]{\rule[0.30em]{#1em}{0.13em}}

% ---- 直撇（审计：44.2 用弯引号、44.3 用直撇，必须统一为直楔撇）----
\newcommand{\FPrime}{\textsf{\textquotesingle}}

% ---- 乘号（审计：46.2 的 × 比源小 25–30%，标签内需按源校准）----
% 用 \FTimes 时在 node 内与底字同字号，并在目检时量宽高。


% ---- 斜杠分数（审计 A：35.1 的 r/3 用 38pt 斜杠手搓）----
\newcommand{\FSlashFrac}[2]{#1\raisebox{-0.18em}{\fontsize{1.35em}{1.35em}\selectfont\textsf{/}}\kern-0.06em #2}

% ---- 交集号 ∩（Heros 无此字，固定回退字号/基线，避免逐图不同）----
\newcommand{\FIntersect}{\raisebox{0.02em}{\fontsize{0.95em}{1em}\selectfont\mfallback{∩}}}

% ---- 括号/开区间（图形层，不是字形）----
% 一律 \draw ... arc[start angle=, end angle=, radius=]：
%   · 端点记号背离它所界的区间；
%   · 左端/右端分别 10× 特写核对；
%   · 闭区间压在轴上时竖脊右移半个线宽；
%   · 同一张图同类记号同一种写法。

% 依赖：fontspec（已 \setsansfont{TeX Gyre Heros}）、graphicx 不需要
% 规则见同目录《符号验收清单.md》
% =====================================================================

% ---- 字体角色（全部走 sans = TeX Gyre Heros）----
\newcommand{\SIt}[1]{\textsf{\textit{#1}}}          % 斜体变量
\newcommand{\SBf}[1]{\textsf{\textbf{#1}}}          % 粗体
\newcommand{\SBfIt}[1]{\textsf{\textbf{\textit{#1}}}} % 粗斜体变量
\newcommand{\SNum}[1]{\textsf{\textup{#1}}}         % 直立数字

% ---- 上下标：字号/位置写死，避免数学样式与手调混用 ----
% 审计结论（2026-10-20）：源图下标/底字高比在 0.77–0.92em，0.72em 偏小。
% 默认 0.80em；仍需按源图逐图标定（44.4/44.5 约 0.85–0.92，39/41/46 约 0.77–0.80）。
\newcommand{\FSup}[2]{#1\raisebox{0.30em}{\fontsize{0.72em}{0.82em}\selectfont #2}}
\newcommand{\FSub}[2]{#1\raisebox{-0.30em}{\fontsize{0.80em}{0.92em}\selectfont #2}}

% ---- 分数：数字直立、方头粗杆；杆宽 = max(分子宽,分母宽) + 2pt ----
% 杆厚 4pt（配合 20–26pt 标签）；小字号图可局部调 0.17em
% 2026-10-20 修正：原 \shortstack 负 kern 不随字号缩放（23pt 下总高仅 30pt、杆贴字），
% 改为 offinterlineskip + em 定间距：上间隙 0.12em、下间隙 0.13em
% （37.2 源图 23pt 实测四标签上/下间隙均值 2.6/4.0pt；取近平衡且总高保持 42–43pt）。
\newcommand{\FFrac}[2]{%
  \begingroup
  \setbox0=\hbox{#1}\setbox2=\hbox{#2}%
  \dimen0=\wd0 \ifdim\wd2>\dimen0 \dimen0=\wd2 \fi
  \advance\dimen0 by 2pt
  \leavevmode\vbox{\offinterlineskip
    \hbox to \dimen0{\hss\box0\hss}%
    \kern0.12em
    \hbox{\rule{\dimen0}{4pt}}%
    \kern0.13em
    \hbox to \dimen0{\hss\box2\hss}%
  }%
  \endgroup}

% ---- 真减号（审计：源图减号 15–18px 宽、3–5px 厚；Heros 常规字重偏细）----
% 可选参数 = 宽度（em）：正文 0.58；上标内用 \FMinus[0.72]
\newcommand{\FMinus}[1][0.58]{\rule[0.30em]{#1em}{0.13em}}

% ---- 直撇（审计：44.2 用弯引号、44.3 用直撇，必须统一为直楔撇）----
\newcommand{\FPrime}{\textsf{\textquotesingle}}

% ---- 乘号（审计：46.2 的 × 比源小 25–30%，标签内需按源校准）----
% 用 \FTimes 时在 node 内与底字同字号，并在目检时量宽高。


% ---- 斜杠分数（审计 A：35.1 的 r/3 用 38pt 斜杠手搓）----
\newcommand{\FSlashFrac}[2]{#1\raisebox{-0.18em}{\fontsize{1.35em}{1.35em}\selectfont\textsf{/}}\kern-0.06em #2}

% ---- 交集号 ∩（Heros 无此字，固定回退字号/基线，避免逐图不同）----
\newcommand{\FIntersect}{\raisebox{0.02em}{\fontsize{0.95em}{1em}\selectfont\mfallback{∩}}}

% ---- 括号/开区间（图形层，不是字形）----
% 一律 \draw ... arc[start angle=, end angle=, radius=]：
%   · 端点记号背离它所界的区间；
%   · 左端/右端分别 10× 特写核对；
%   · 闭区间压在轴上时竖脊右移半个线宽；
%   · 同一张图同类记号同一种写法。

% 依赖：fontspec（已 \setsansfont{TeX Gyre Heros}）、graphicx 不需要
% 规则见同目录《符号验收清单.md》
% =====================================================================

% ---- 字体角色（全部走 sans = TeX Gyre Heros）----
\newcommand{\SIt}[1]{\textsf{\textit{#1}}}          % 斜体变量
\newcommand{\SBf}[1]{\textsf{\textbf{#1}}}          % 粗体
\newcommand{\SBfIt}[1]{\textsf{\textbf{\textit{#1}}}} % 粗斜体变量
\newcommand{\SNum}[1]{\textsf{\textup{#1}}}         % 直立数字

% ---- 上下标：字号/位置写死，避免数学样式与手调混用 ----
% 审计结论（2026-10-20）：源图下标/底字高比在 0.77–0.92em，0.72em 偏小。
% 默认 0.80em；仍需按源图逐图标定（44.4/44.5 约 0.85–0.92，39/41/46 约 0.77–0.80）。
\newcommand{\FSup}[2]{#1\raisebox{0.30em}{\fontsize{0.72em}{0.82em}\selectfont #2}}
\newcommand{\FSub}[2]{#1\raisebox{-0.30em}{\fontsize{0.80em}{0.92em}\selectfont #2}}

% ---- 分数：数字直立、方头粗杆；杆宽 = max(分子宽,分母宽) + 2pt ----
% 杆厚 4pt（配合 20–26pt 标签）；小字号图可局部调 0.17em
% 2026-10-20 修正：原 \shortstack 负 kern 不随字号缩放（23pt 下总高仅 30pt、杆贴字），
% 改为 offinterlineskip + em 定间距：上间隙 0.12em、下间隙 0.13em
% （37.2 源图 23pt 实测四标签上/下间隙均值 2.6/4.0pt；取近平衡且总高保持 42–43pt）。
\newcommand{\FFrac}[2]{%
  \begingroup
  \setbox0=\hbox{#1}\setbox2=\hbox{#2}%
  \dimen0=\wd0 \ifdim\wd2>\dimen0 \dimen0=\wd2 \fi
  \advance\dimen0 by 2pt
  \leavevmode\vbox{\offinterlineskip
    \hbox to \dimen0{\hss\box0\hss}%
    \kern0.12em
    \hbox{\rule{\dimen0}{4pt}}%
    \kern0.13em
    \hbox to \dimen0{\hss\box2\hss}%
  }%
  \endgroup}

% ---- 真减号（审计：源图减号 15–18px 宽、3–5px 厚；Heros 常规字重偏细）----
% 可选参数 = 宽度（em）：正文 0.58；上标内用 \FMinus[0.72]
\newcommand{\FMinus}[1][0.58]{\rule[0.30em]{#1em}{0.13em}}

% ---- 直撇（审计：44.2 用弯引号、44.3 用直撇，必须统一为直楔撇）----
\newcommand{\FPrime}{\textsf{\textquotesingle}}

% ---- 乘号（审计：46.2 的 × 比源小 25–30%，标签内需按源校准）----
% 用 \FTimes 时在 node 内与底字同字号，并在目检时量宽高。


% ---- 斜杠分数（审计 A：35.1 的 r/3 用 38pt 斜杠手搓）----
\newcommand{\FSlashFrac}[2]{#1\raisebox{-0.18em}{\fontsize{1.35em}{1.35em}\selectfont\textsf{/}}\kern-0.06em #2}

% ---- 交集号 ∩（Heros 无此字，固定回退字号/基线，避免逐图不同）----
\newcommand{\FIntersect}{\raisebox{0.02em}{\fontsize{0.95em}{1em}\selectfont\mfallback{∩}}}

% ---- 括号/开区间（图形层，不是字形）----
% 一律 \draw ... arc[start angle=, end angle=, radius=]：
%   · 端点记号背离它所界的区间；
%   · 左端/右端分别 10× 特写核对；
%   · 闭区间压在轴上时竖脊右移半个线宽；
%   · 同一张图同类记号同一种写法。


% 三点折线：V / Λ
\newcommand{\tri}[6]{\draw (#1,#2) -- (#3,#4) -- (#5,#6);}
% 贴竖格线的折角：格线 x、上下端 y、顶点 (vx,vy)
\newcommand{\chev}[5]{\draw (#1,#2) -- (#4,#5) -- (#1,#3);}

\begin{document}
\begin{tikzpicture}[x=1pt, y=-1pt, line width=5pt, line cap=round, line join=round]

  %% ===== 1. 标签（所有几何之前，白底）=====
  \node[anchor=center, inner sep=0pt, fill=white] at (565,113)
    {\fontsize{26}{31}\selectfont\FSub{\SBfIt f}{\SBf{\SNum 3}}};

  %% ===== 2. 刻度梳：65 齿，每第 4 根长 =====
  \draw (27,236) -- (490,236);
  \foreach \i in {0,...,64} {
    \pgfmathsetmacro\x{28.5+7.1953*\i}
    \ifnum\numexpr\i-4*(\i/4)=0
      \draw (\x,236) -- (\x,199.5);
    \else
      \draw (\x,236) -- (\x,224.5);
    \fi
  }

  %% ===== 3. 弧形箭头（后掠头，4pt）=====
  \draw[-{Stealth[length=14pt,width=10pt]}, line width=4pt]
    (493,167) .. controls (530,144) and (586,134) .. (637,159);

  %% ===== 4. 4x4 方格骨架 =====
  % 4x4 网格等分路径所在正方形；线位与 Peano 路径的顶点取同一组坐标
  \draw (670.5,10) rectangle (1106.5,452.5);
  \foreach \x in {779.5,888.5,997.5} \draw (\x,10) -- (\x,452.5);
  \foreach \y in {120.6,231.2,341.9} \draw (670.5,\y) -- (1106.5,\y);

  %% ===== 5. Peano 曲线 f3：一条连续路径（64 条小三角路径首尾相接）=====
  \draw (670.5,452.5) -- (697.8,424.8) -- (670.5,397.2) -- (697.8,369.5) -- (725.0,397.2) -- (752.2,369.5) -- (779.5,397.2) -- (752.2,424.8) -- (779.5,452.5) -- (806.8,424.8) -- (834.0,452.5) -- (861.2,424.8) -- (834.0,397.2) -- (861.2,369.5) -- (834.0,341.9) -- (806.8,369.5) -- (779.5,341.9) -- (806.8,314.2) -- (834.0,341.9) -- (861.2,314.2) -- (834.0,286.6) -- (861.2,258.9) -- (834.0,231.2) -- (806.8,258.9) -- (779.5,231.2) -- (752.2,258.9) -- (779.5,286.6) -- (752.2,314.2) -- (725.0,286.6) -- (697.8,314.2) -- (670.5,286.6) -- (697.8,258.9) -- (670.5,231.2) -- (697.8,203.6) -- (725.0,231.2) -- (752.2,203.6) -- (725.0,175.9) -- (752.2,148.3) -- (725.0,120.6) -- (697.8,148.3) -- (670.5,120.6) -- (697.8,93.0) -- (670.5,65.3) -- (697.8,37.7) -- (725.0,65.3) -- (752.2,37.7) -- (779.5,65.3) -- (752.2,93.0) -- (779.5,120.6) -- (806.8,93.0) -- (779.5,65.3) -- (806.8,37.7) -- (834.0,65.3) -- (861.2,37.7) -- (888.5,65.3) -- (861.2,93.0) -- (888.5,120.6) -- (861.2,148.3) -- (834.0,120.6) -- (806.8,148.3) -- (834.0,175.9) -- (806.8,203.6) -- (834.0,231.2) -- (861.2,203.6) -- (888.5,231.2) -- (915.8,203.6) -- (943.0,231.2) -- (970.2,203.6) -- (943.0,175.9) -- (970.2,148.3) -- (943.0,120.6) -- (915.8,148.3) -- (888.5,120.6) -- (915.8,93.0) -- (888.5,65.3) -- (915.8,37.7) -- (943.0,65.3) -- (970.2,37.7) -- (997.5,65.3) -- (970.2,93.0) -- (997.5,120.6) -- (1024.8,93.0) -- (997.5,65.3) -- (1024.8,37.7) -- (1052.0,65.3) -- (1079.2,37.7) -- (1106.5,65.3) -- (1079.2,93.0) -- (1106.5,120.6) -- (1079.2,148.3) -- (1052.0,120.6) -- (1024.8,148.3) -- (1052.0,175.9) -- (1024.8,203.6) -- (1052.0,231.2) -- (1079.2,203.6) -- (1106.5,231.2) -- (1079.2,258.9) -- (1106.5,286.6) -- (1079.2,314.2) -- (1052.0,286.6) -- (1024.8,314.2) -- (997.5,286.6) -- (1024.8,258.9) -- (997.5,231.2) -- (970.2,258.9) -- (943.0,231.2) -- (915.8,258.9) -- (943.0,286.6) -- (915.8,314.2) -- (943.0,341.9) -- (970.2,314.2) -- (997.5,341.9) -- (970.2,369.5) -- (943.0,341.9) -- (915.8,369.5) -- (943.0,397.2) -- (915.8,424.8) -- (943.0,452.5) -- (970.2,424.8) -- (997.5,452.5) -- (1024.8,424.8) -- (997.5,397.2) -- (1024.8,369.5) -- (1052.0,397.2) -- (1079.2,369.5) -- (1106.5,397.2) -- (1079.2,424.8) -- (1106.5,452.5);

  %% ===== 钉死画布 =====
  \pgfresetboundingbox
  \path[use as bounding box] (0,0) rectangle (1135,459);
\end{tikzpicture}
\end{document}
